\section{Automorphisms of the rank one Hodge moduli space}\label{auto-hodge}

\subsection{Connected component of the automorphism group}

Let $X$ be an irreducible smooth complex projective curve of genus $g$, with $g \geq 1$. The holomorphic cotangent bundle of $X$ will be denoted by $K_X$.
\begin{Definition} For any $\lambda \in \mathbb C$, a rank one $\lambda$-connection on $X$ is a pair of the form $(L, D)$, where $L$ is a holomorphic line bundle on~$X$ and
\begin{gather*}
D \colon \ L \longrightarrow L\otimes K_X
\end{gather*}
is a holomorphic dif\/ferential operator of order no more than one such that
\begin{gather*}
D(f_0s) = f_0D(s)+\lambda\cdot s\otimes (df_0)
\end{gather*}
for every locally def\/ined holomorphic section $s$ of $L$ and every locally def\/ined holomorphic function $f_0$ on $X$.
\end{Definition}

So, if $(L, D)$ is a $\lambda$-connection with $\lambda \not= 0$, then $D/\lambda$ is a holomorphic connection on $L$; a~$0$-connection $D$ on $L$ is an element of $H^0(X, K_X)$ because such a $D$ is ${\mathcal O}_X$-linear. Note that for a~$\lambda$-connection $(L, D)$, with $\lambda \not= 0$ we have $\text{degree}(L) = 0$ because $L$ admits a holomorphic connection. Also, note that any holomorphic connection on a Riemann surface is automatically f\/lat.

If $(L, D)$ is a $0$-connection we will always impose the condition that $\text{degree}(L) = 0$.

Let
\begin{gather*}
{\mathcal M}_{\rm Hod} = {\mathcal M}^X_{\rm Hod}(1)
\end{gather*}
be the Hodge moduli space of rank one $\lambda$-connections. So, the points of ${\mathcal M}_{\rm Hod}$ parametrize triples of the form $(\lambda, L, D)$, where $\lambda \in \mathbb C$, $L$ is a holomorphic line bundle on $X$ of degree zero and~$D$ is a $\lambda$-connection on~$L$.

Def\/ine an algebraic morphism
\begin{gather}\label{lf}
f \colon \ {\mathcal M}_{\rm Hod} \longrightarrow {\mathbb C} , \qquad (\lambda, L, D) \longmapsto \lambda .
\end{gather}
For any $\lambda \in \mathbb C$, let
\begin{gather*}%\label{fiber}
f^\lambda := f^{-1}(\lambda) \subset {\mathcal M}_{\rm Hod}
\end{gather*}
be the f\/iber over $\lambda$. By def\/inition, the f\/iber $f^0$ is the moduli space of Higgs line bundles of degree zero $\operatorname{Pic}^0(X)\times H^0(X, K_X)$ on $X$, and the f\/iber $f^1$ is the moduli space of rank one holomorphic connections on $X$. The latter space is also called the de Rham moduli space. For $\lambda \not= 0$, the f\/iber $f^\lambda$ is canonically identif\/ied with $f^1$ by the map $(L, D) \longmapsto (L, D/\lambda)$. We note that $f^1$ is biholomorphic to the Betti moduli space $\operatorname{Hom}(\pi_1(X), {\mathbb C}^*) = ({\mathbb C}^*)^{2g}$ by sending a holomorphic connection to its monodromy representation, and hence $f^1$ does not admit a~compact submanifold of positive dimension. Since $\operatorname{Pic}^0(X)\times H^0(X, K_X)$ has the projective variety $\operatorname{Pic}^0(X)$ of positive dimension as a~subvariety, it follows that~$f^0$ is not biholomorphic to~$f^1$.

The group of all algebraic automorphisms of the quasiprojective variety
${\mathcal M}_{\rm Hod}$ will be denoted by $\operatorname{Aut}({\mathcal M}_{\rm Hod})$. Let
\begin{gather*}
\operatorname{Aut}({\mathcal M}_{\rm Hod})_0 \subset \operatorname{Aut}({\mathcal M}_{\rm Hod})
\end{gather*}
be the connected component containing the identity automorphism and
 let
\begin{gather*}
{\mathbb V} := \operatorname{Morphisms}\big({\mathbb C}, H^0(X, K_X)\big)
\end{gather*}
be the inf\/inite dimensional complex vector space parametrizing all algebraic maps from $\mathbb C$ to the af\/f\/ine space $H^0(X, K_X)$. Then we have the following theorem.
\begin{Theorem}\label{thm1}
The group $\operatorname{Aut}({\mathcal M}_{\rm Hod})_0$ fits in a short exact sequence of groups
\begin{gather*}
0 \longrightarrow {\mathbb V} \stackrel{\iota}{\longrightarrow} \operatorname{Aut}({\mathcal M}_{\rm Hod})_0 \stackrel{h}{\longrightarrow}
\operatorname{Pic}^0(X)\times {\mathbb C}^* \longrightarrow 0 .
\end{gather*}
\end{Theorem}

\begin{proof} Let
\begin{gather*}
T \colon \ {\mathcal M}_{\rm Hod} \longrightarrow {\mathcal M}_{\rm Hod}
\end{gather*}
be any algebraic automorphism that lies in $\operatorname{Aut}({\mathcal M}_{\rm Hod})_0$. It is known that the f\/iber $f^1$ does not admit any nonconstant algebraic function \cite[Proposition~2.2]{BBS}. Since $f^\lambda$ is isomorphic to $f^1$ for every $\lambda \in {\mathbb C}^*$, it follows that $f^\lambda$ does not admit any nonconstant algebraic function if $\lambda \in {\mathbb C}^*$. Hence the composition
\begin{gather}\label{e1}
f^\lambda \stackrel{T\vert_{f^\lambda}}{\longrightarrow} {\mathcal M}_{\rm Hod} \stackrel{f}{\longrightarrow} \mathbb C
\end{gather}
is a constant function for all $\lambda \in {\mathbb C}^*$. Note that the point $f\circ T(f^{\lambda})$ lies in ${\mathbb C}^*$ because $f^0$ is not isomorphic to $f^\lambda$. Since this also holds for~$T^{-1}$, it follows that the composition in~\eqref{e1} is a~constant function for each $\lambda \in {\mathbb C}$. This implies that there is an algebraic automorphism
\begin{gather*}
\tau_0 \colon \ {\mathbb C} \longrightarrow \mathbb C
\end{gather*}
such that
\begin{gather*}
f\circ T = \tau_0 \circ f .
\end{gather*}
Since $f^0$ is not isomorphic to any $f^\lambda$, $\lambda \in {\mathbb C}^*$, it follows that $\tau_0(0) = 0$. Therefore, $\tau_0$ coincides with the multiplication of $\mathbb C$ by a f\/ixed nonzero number; let
\begin{gather}\label{e4}
\tau \in {\mathbb C}^*
\end{gather}
be such that $\tau_0(z) = \tau \cdot z$ for all $z \in \mathbb C$.

The group of all automorphisms of the variety $\operatorname{Pic}^0(X)$ will be denoted by $\operatorname{Aut}(\operatorname{Pic}^0(X))$ (since the variety $\operatorname{Pic}^0(X)$ is projective, any holomorphic automorphism of it is algebraic). The connected component of~$\operatorname{Aut}(\operatorname{Pic}^0(X))$ containing the identity automorphism will be denoted by~$\operatorname{Aut}(\operatorname{Pic}^0(X))_0$. The automorphisms of $\operatorname{Pic}^0(X)$ given by translations of the group $\operatorname{Pic}^0(X)$ lie in $\operatorname{Aut}(\operatorname{Pic}^0(X))_0$. In fact, this way $\operatorname{Pic}^0(X)$ gets identif\/ied with
$\operatorname{Aut}(\operatorname{Pic}^0(X))_0$; note that the Lie algebra $H^0(\operatorname{Pic}^0(X), T\operatorname{Pic}^0(X))$ of~$\operatorname{Aut}(\operatorname{Pic}^0(X))_0$ coincides with the Lie algebra of $\operatorname{Pic}^0(X)$ by evaluating sections of~$T\operatorname{Pic}^0(X)$ at the identity element.

For any $\lambda \in \mathbb C$, let $\phi^\lambda$ be the morphism def\/ined by
\begin{gather}\label{e7}
\phi^\lambda \colon \ f^\lambda \longrightarrow \operatorname{Pic}^0(X) ,\qquad (L, D) \longmapsto L .
\end{gather}
The f\/ibers of $\phi^\lambda$ are isomorphic to $H^0(X, K_X)$, because the space of all holomorphic connections on a line bundle $L \in \operatorname{Pic}^0(X)$ is an af\/f\/ine space for the vector space $H^0(X, K_X)$, and $H^0(X, K_X)$ is the space of all Higgs f\/ields on any line bundle. There is no nonconstant algebraic map from an af\/f\/ine space to an abelian variety. Hence, there is no nonconstant algebraic map from a f\/iber of $\phi^\lambda$ to $\operatorname{Pic}^0(X)$. Consequently, we get a map
\begin{gather*}
\Phi \colon \  {\mathbb C} \longrightarrow \operatorname{Aut}\big(\operatorname{Pic}^0(X)\big),
\end{gather*}
which is uniquely determined by the condition that the following diagram is commutative
\begin{gather*}
\begin{matrix}
f^\lambda & \stackrel{T\vert_{f^\lambda}}{\longrightarrow} & f^{\tau\lambda}\\
~ ~~ \Big\downarrow \phi^\lambda && ~ ~~~~ \Big\downarrow \phi^{\tau\lambda}\\
\operatorname{Pic}^0(X) & \stackrel{\Phi(\lambda)}{\longrightarrow} &\operatorname{Pic}^0(X)
\end{matrix}
\end{gather*}
for all $\lambda \in \mathbb C$, where $\tau$ is the complex number in \eqref{e4}$ $. Note that since $T \in \operatorname{Aut}({\mathcal M}_{\rm Hod})_0$, it follows that the image of $\Phi$ lies in $\operatorname{Aut}\big(\operatorname{Pic}^0(X)\big)_0 = \operatorname{Pic}^0(X) \subset \operatorname{Aut}\big(\operatorname{Pic}^0(X)\big)$. Thus, $\Phi$ is a~constant map, again because there is no nonconstant algebraic map from $\mathbb C$ to $\operatorname{Pic}^0(X)$. Let{\samepage
\begin{gather}\label{e5}
\widehat{\Phi} = \Phi({\mathbb C}) \in \operatorname{Pic}^0(X)
\end{gather}
be the image of the map $\Phi$.}

Let
\begin{gather*}
h \colon \ \operatorname{Aut}({\mathcal M}_{\rm Hod})_0 \longrightarrow \operatorname{Pic}^0(X)\times {\mathbb C}^*
\end{gather*}
be the map def\/ined by $T \longmapsto \big(\widehat{\Phi}, \tau\big)$. It is straight-forward to check that $h$ is a group homomorphism.

We will now show that $h$ is surjective. For this, f\/irst note that the multiplicative action of~${\mathbb C}^*$ on $\mathbb C$ has a natural lift to an action of ${\mathbb C}^*$ on ${\mathcal M}_{\rm Hod}$. Indeed, any $c \in {\mathbb C}^*$ acts on ${\mathcal M}_{\rm Hod}$ as
\begin{gather*}
(\lambda, L, D) \longmapsto (c\cdot\lambda, L, c\cdot D) .
\end{gather*}
For all these automorphisms of ${\mathcal M}_{\rm Hod}$ given by the action of ${\mathbb C}^*$, the corresponding elements of~$\operatorname{Pic}^0(X)$ (see \eqref{e5}) coincide with the identity element. Therefore, to prove that~$h$ is surjective, it suf\/f\/ices to show that the composition of $h$ with the projection $\operatorname{Pic}^0(X)\times {\mathbb C}^*  \longrightarrow \operatorname{Pic}^0(X)$ is surjective. Take any $L_0 \in \operatorname{Pic}^0(X)$ and f\/ix a holomorphic connection~$D_0$ on~$L_0$. Def\/ine
\begin{gather*}
\beta \colon \  {\mathcal M}_{\rm Hod} \longrightarrow {\mathcal M}_{\rm Hod} ,\qquad
(\lambda, L, D) \longmapsto (\lambda, L\otimes L_0, (D\otimes \text{Id}_{L_0}) +(\lambda\cdot \text{Id}_L\otimes D_0)) .
\end{gather*}
It is straight-forward to check that $\beta \in \operatorname{Aut}({\mathcal M}_{\rm Hod})$; its inverse is the corresponding map for $(L_0^*, D_0^*)$. Since the moduli space of rank one connections on $X$ is connected, it follows that $\beta \in \operatorname{Aut}({\mathcal M}_{\rm Hod})_0$; note that $\beta$ for the trivial connection $({\mathcal O}_X, d)$ is the identity map of ${\mathcal M}_{\rm Hod}$. The element in $\operatorname{Pic}^0(X)$ corresponding to $\beta$ (see~\eqref{e5}) is clearly $L_0$. Therefore, the composition of $h$ with the projection $\operatorname{Pic}^0(X)\times {\mathbb C}^*  \longrightarrow \operatorname{Pic}^0(X)$ and hence $h$ itself are surjective.

Next, we need to def\/ine ${\mathbb V} \stackrel{\iota}{\longrightarrow} \operatorname{Aut}({\mathcal M}_{\rm Hod})_0$. To do so, consider for $v \in \mathbb V$ the automorphism
\begin{gather}\label{e6}
\iota_v \colon \ {\mathcal M}_{\rm Hod} \longrightarrow {\mathcal M}_{\rm Hod} ,\qquad (\lambda, L, D) \longmapsto (\lambda, L, D+v(\lambda)) .
\end{gather}
Clearly, we have $\iota_v \in \operatorname{Aut}({\mathcal M}_{\rm Hod})_0$, and also $\iota_v \in \operatorname{kernel}(h)$. Therefore, there is an injective homomorphism
\begin{gather*}
\iota \colon \ {\mathbb V} \longrightarrow \operatorname{kernel}(h) \subset \operatorname{Aut}({\mathcal M}_{\rm Hod})_0 ,\qquad v \longmapsto \iota_v .
\end{gather*}
We will prove that $\operatorname{image}(\iota) = \operatorname{kernel}(h)$.

In order to prove that $\operatorname{image}(\iota) = \operatorname{kernel}(h)$, take any $T \in \operatorname{Aut}({\mathcal M}_{\rm Hod})_0$ such that
\begin{gather*}T \in \operatorname{kernel}(h) .\end{gather*} We have $T(f^\lambda) = f^\lambda$ for all
$\lambda \in {\mathbb C}^*$ because the constant in \eqref{e4} for $T$ is~$1$. Since the element in \eqref{e5} corresponding to $T$ is the trivial line bundle, we get a morphism
\begin{gather*}
\delta_\lambda \colon \  f^\lambda \longrightarrow H^0(X, K_X) ,\qquad z \longmapsto T(z) - z ;
\end{gather*}
note that since the f\/ibers of $\phi^\lambda$ (constructed in \eqref{e7}) are af\/f\/ine spaces for $H^0(X, K_X)$, and $\phi^\lambda(T(z)) = \phi^\lambda(z)$ for all $z \in f^\lambda$, we have $T(z) - z \in H^0(X, K_X)$. As there are no nonconstant algebraic functions on $f^\lambda$ \cite[Proposition~2.2]{BBS}, it follows that the above function $\delta_\lambda$ is a constant one.

All automorphism of the af\/f\/ine space $H^0(X, K_X)$ are of the form $u \longmapsto A(u)+u_0$, where $A \in \text{GL}(H^0(X, K_X))$ and $u_0 \in H^0(X, K_X)$. Considering the restrictions of $T$ to open subsets of the form $h^{-1}(V\times U)$, where $U$ is an analytic neighborhood of $0 \in \mathbb C$ and $V$ is an analytic open subset of $\operatorname{Pic}^0(X)$, it now follows that the above map
\begin{gather*}\lambda \longmapsto \operatorname{image}(\delta_\lambda) \in H^0(X, K_X)\end{gather*}
extends across $0 \in \mathbb C$ as a holomorphic map from $\mathbb C$ to $H^0(X, K_X)$. In other words, there is a~unique element $v \in {\mathbb V}$ such that $v(\lambda) = \operatorname{image} (\delta_\lambda)$ for all $\lambda \in {\mathbb C}^*$. Clearly, we have $\iota_v = T$, where~$\iota_v$ is constructed in~\eqref{e6}. This implies that $\operatorname{image}(\iota) = \operatorname{kernel}(h)$. This completes the proof of the theorem.
\end{proof}
